%%% Jim Ashe (jim.ashe@wgu.edu; x4209)%%%%
\documentclass{exam}
\usepackage{WGU_worksheet_preamble}
\usepackage{footnote}
\usepackage[symbol]{footmisc}

%%header%%%%%%%%%%%%%%%%%%%%%%%%%%%%%%%%%%
\headrule
\header{ 
    \href{mailto:jim.ashe@wgu.edu;cemal.tepe@wgu.edu?subject=C949: Big-O worksheet}{Dr.\,Jim Ashe \& Dr.\,Cemal Tepe}
    \\ \textbf{C949 Data Structures \& Algorithms}}
{Big-$\mathcal{O}$ Worksheet}
{}
%%%%%%%%%%%%%%%%%%%%%%%%%%%%%%%%%%%%%%%%%
\lfoot{\footnotesize Suggestions or errors? Contact: \href{mailto:jim.ashe@wgu.edu?subject=C949: Big-O worksheet}{Dr.\,Jim Ashe}.}
\cfoot[]{Page \thepage\ of \numpages}
\rfoot{\today}
%%%%%%%%%%%%%%%%%%%%%%%%%%%%%%%%%%%%%%%%%
\newcommand{\itemSpace}{36pt} %adjusts spacing between questions

% \printanswers
\shadedsolutions
\colorlet{SolutionColor}{lightBlueWGU!15!white}

\begin{document}

\subsection*{Basic Properties of Big-$\mathcal{O}$}
    \vspace{-.1cm}
    \begin{multicols}{2}
        Time complexity describes the amount of computational time an algorithm takes. As this can vary and run times are fast for small $n$, we describe the behavior as $n$ gets large for best, average, and worst-case scenarios -mainly, we focus on worst, i.e., Big-$O$. Here's a list of some functions listed from fast to slow\footnote[1]{You don't need to know these for the C949 assessment! But any computer science major should be familiar with them.}:
        \begin{itemize}
            \item[] $\mathcal{O}(1)$ \textbf{Constant}
            \item[] $\mathcal{O}(\log(n))$ \textbf{Logarithmic}. 
            \begin{itemize}
                \item[] Ignore Log bases: $\log_{10}(n)=\log(n)$
                \item[] Same growth as $\log(n^{k})=k\log(n)$ 
            \end{itemize}
            
            \item[] $\mathcal{O}(\log^{c}(n))$ \textbf{Polylogarithmic}
            \begin{itemize}
                \item[] Don't forget that: $\log^{c}x=(\log x)^{c}$.
            \end{itemize}
            \item[] $\mathcal{O}(n)$ \textbf{Linear}
            \item[] $\mathcal{O}(n\log(n))$ \textbf{Linearithmic/Loglinear}
            \item[] $\mathcal{O}(n^{2})$ \textbf{Quadratic}
            \item[] $\mathcal{O}(n^{2}\log(n))$
            \item[] $\mathcal{O}(n^{3})$ \textbf{Polynomial}
            \item[] $\mathcal{O}(2^{n})$ \textbf{Exponential} 
            \item[] $\mathcal{O}(n!)$ \textbf{Factorial}
        \end{itemize}
        \columnbreak
        \begin{tikzpicture} 
            \begin{axis}[height=6cm, 
            	width=6cm,scale=1.0,
            	axis lines=middle,
            	xticklabels={,,}, yticklabels={,,},
            	axis line style={-latex}, 
            	xlabel={$n\rightarrow \infty$},
            	ylabel={$f(n)$},             	
            	x label style={at={(axis description cs:0.5,-0.05)},anchor=north},
            	y label style={at={(axis description cs:-0.125,.5)},anchor=south},
            	clip=false, %doesn't clip outside graph
            	%domain=0:100,
            	ymin=0,xmin=0,
            	ymax = 150,xmax=150
            	]
            	\addplot[teal, thick, samples=10, domain=0:18.835]{1.3^x} node[above]{$c^{n}$};
            	\addplot[Magenta, thick, domain=0:52.1936]{.25*x^1.6} node[above]{$c\cdot n^{10}$};
            	\addplot[blue, thick, domain=0:91.8315]{.25*x^1.4} node[right]{$c\cdot n^{2}$};
            	\addplot[RedOrange, thick, domain=0:140]{.13*(x+1)*ln(x+1)/ln(2)} node[right]{$n\log_{c}n$};
            	\addplot[red, thick, domain=0:140] {.5*x} node[right]{$c\cdot n$};
            	\addplot[purple, thick, domain=0:140]{3.5*ln(x)/ln(2)} node[right]{$\log_{c}n$};
            	\addplot[OliveGreen, thick, domain=0:140]{10} node[right]{$c$};
            \end{axis}
        \end{tikzpicture}
        \\ \textbf{Some properties} \\
            We ignore coefficients and log bases. 
                \begin{itemize}
                    \item [\textbf{ex.}] $7x^{2}$ is $\mathcal{O}(x^{2})$ and $3\log_{20}{x}$ is $\mathcal{O}(\log{x})$ 
                        \footnote[7]{\textbf{Clarification of notation: }The following statements all mean the same thing:
                            \begin{itemize}
                                \item[] “$f(x)$ is $\mathcal{O}(g(x))$” OR  
                                \item[] “$f(x)$ is of $\mathcal{O}(g(x))$” OR 
                                \item[] “$f(x)=\mathcal{O}(g(x))$” OR
                                \item[] “$f(x)\in \mathcal{O}(g(x))$” 
                            \end{itemize}
                        $\mathcal{O}(g(x))$ is a collection of functions (i.e. a set) so we should write "$f(x)\in \mathcal{O}(g(x))$", but "$f(x)=\mathcal{O}(g(x))$” is commonly used and what you'll see on the assessment. This can be confusing since $\mathcal{O}(n)=\mathcal{O}(n^{2})$ but $\mathcal{O}(n^{2})\neq \mathcal{O}(n)$!! See this interesting thread on StackExchange \href{https://math.stackexchange.com/questions/2066004/big-o-notation-is-element-of-or-is-equal}{Big O Notation “is element of” or “is equal”}. Note the Wiki cites \href{https://en.wikipedia.org/wiki/Donald_Knuth}{Donald Knuth}.
                        } 
                \end{itemize}
        
            When adding functions, we only take the slowest function's Big-$\mathcal{O}$.
            \begin{itemize}
                \item[\textbf{ex.}] $0.01x^3+12x^{2}-2\log{37x}$ is $\mathcal{O}(n^{3})$
            \end{itemize}
            
            When multiplying functions, we take take the product of the functions' Big-$\mathcal{O}$. 
            \begin{itemize}
                \item[\textbf{ex.}] $x^{2}+2\in \mathcal{O}(n^{2})$ and $\log_{10}(x)\in \mathcal{O}(\log{n})$ so then \\
                $(x^{2}+2)(\log_{10}(x)) \text{ is } \mathcal{O}(n^{2} \log(n))$
            \end{itemize}
    \end{multicols}
    
\vspace{-.4cm}
%%use \footnotemark and footnote[n]{} if not in myThm enviro. 
\begin{figure}[H]
    \begin{myThm}
        \centering
            \rowcolors{2}{white}{lightBlueWGU!25}
            % \resizebox{.5\textwidth}{!}{
            \begin{tabular}{cccc}
                \rowcolor{blueWGU!50}
                \textcolor{white}{Sort function} & 
                \textcolor{white}{Average/Typical} &
                \textcolor{white}{Big-$\mathcal{O}$/worst} &
                \textcolor{white}{Fast?}\\
                Bubble sort & $\Theta(n^{2})$ & $\mathcal{O}(n^{2})$ & No\\
                Selection sort & $\Theta(n^{2})$ & $\mathcal{O}(n^{2})$ & No\\
                Insertion sort & $\Theta(n^{2})$ & $\mathcal{O}(n^{2})$ & No\\
                Quick sort & $\Theta(n\log(n))$ & \textcolor{red}{$\mathcal{O}(n^{2})$} & Yes\\
                Bucket sort & $\Theta(n)$ & \textcolor{red}{$\mathcal{O}(n^{2})$} & Yes\\
                Heap sort & $\Theta(n\log(n))$ & $\mathcal{O}(n\log(n))$ & Yes\\
                Merge sort & $\Theta(n\log(n))$ & $\mathcal{O}(n\log(n))$ & Yes\\
                Radix sort\footnote{Actually, as shown in most examples Radix is $\mathcal{O}(\left \lfloor\log_{b}(k) \right \rfloor (n+b))$ where $b$ is the base (or bucket size) and $k$ is the max possible value, but this can be reduced to $\mathcal{O}(n)$ by increasing $b$ to $n$ -\emph{use $\mathcal{O}(n)$ for the assessments}.} & $\Theta(n)$ & $\mathcal{O}(n)$ & Yes 
            \end{tabular}
            % }
        \caption*{Sorting algorithms time complexities to memorize.\footnote{The assessments ask about both worst and average time complexities. They also mistakenly use Big-$\mathcal{O}$ notation for both average and worst time complexities. However, for this list, they  only differ for \textit{Quicksort} and \textit{Bucket}.}}
    \end{myThm}
\end{figure}

\begin{figure}
    \begin{myThm}
        \centering
            \rowcolors{2}{white}{lightBlueWGU!25}
            % \resizebox{.5\textwidth}{!}{
            \begin{tabular}{cc}
                \rowcolor{blueWGU!50}
                \textcolor{white}{Sort function} & \textcolor{white}{Characteristic}\\
                Bubble sort & values are swapped\\
                Bucket sort & values are distributed into sets or "buckets"\\
                Quicksort & look for keywords "pivot"\\
                Merge sort & the list is continually "split"\\
                Radix sort & sorting by least/most significant digits
            \end{tabular}
            % }
        \caption*{Sorting algorithms characteristics \\to identify from code}
    \end{myThm}
\end{figure}

% \footnotetext[4]{The pre-assessment and presumably the assessment ask about both worst and average time complexities. They also mistakenly use Big-$\mathcal{O}$ notation for both average and worst time complexities. However, for this list they  only differ for \textit{Quicksort} and \textit{Bucket}.}

% \footnotetext[3]{Actually Radix is $\Theta(n)$ and $\mathcal{O}(n)$ where $k$ is dependent on the number of bits required to store each key, but the assessment will not make this distinction.}

\newpage
\subsection*{Big-$\mathcal{O}$ problems}

\begin{myThm}
    On the assessment, you will need to determine the Big-$\mathcal{O}$ time-complexity from psuedocode. \footnote{Pseudocode on the assessment won't be written as nicely as it is here! Furthermore, it may be given in the form of Java, C, or Python.} Answers seem to mainly be limited to $\mathcal{O}(1)$, $\mathcal{O}(n)$, $\mathcal{O}(n^{2})$, and $\mathcal{O}(n^3)$. A nutshell summary for the assessment:
    \begin{itemize}
        \item If the number of steps stays the same no matter how large, then it’s constant time: $\mathcal{O}(1)$
        \item If you go through a $n$ long list (linearly; say a loop), then you’ll take at most $n$ steps: $\mathcal{O}(n)$.
        \item A loop in a loop will be $\mathcal{O}(n^{2})$. A loop in a loop in a loop will be $\mathcal{O}(n^{3})$. Often you simply need to count the number of \texttt{for} and \texttt{while} loops.
    \end{itemize}
\end{myThm}

\begin{questions}
\setlength\itemsep{\itemSpace} %set length above question. Needs to be set for each question environment

\question Using the provided pseudo-code, find the worst case performance in Big-$\mathcal{O}$ notation.
        \begin{algorithm}[H]
        \caption{Some messy pseudo-code}\label{algor1}
            \begin{algorithmic}[1]
                \Procedure{SomeProcedure}{}
                    \For{i=1 and 1$<$=n}
                        \State j=1
                        \While{j$<$n}
                            \State j=j+2
                        \EndWhile
                    \EndFor
                \EndProcedure
            \end{algorithmic}
        \end{algorithm}
    \begin{oneparchoices}
        \choice $\mathcal{O}(\log n)$
        \choice $\mathcal{O}(n\log n)$
        \CorrectChoice $\mathcal{O}(n^{2})$
        \choice $\mathcal{O}(n)$
    \end{oneparchoices}
    
\question Using the provided pseudo-code, find the worst case performance in Big-$\mathcal{O}$ notation.
        \begin{algorithm}[H]
        \caption{Some more messy pseudocode}\label{algor2}
            \begin{algorithmic}[1]
                \Procedure{PrintHello}{}
                    \State n=100,000,000
                    \For{i=0; i$<$n}
                        \State Output "Hello"
                        \State i=i+1
                    \EndFor
                \EndProcedure
            \end{algorithmic}
        \end{algorithm}
    \begin{oneparchoices}
        \choice $\mathcal{O}(100,000,000)$
        \choice $\mathcal{O}(n\log n)$
        \choice $\mathcal{O}(n^{2})$
        \CorrectChoice $\mathcal{O}(1)$
    \end{oneparchoices}

\question Using the provided pseudo-code, find the worst case performance in Big-$\mathcal{O}$ notation.
        \begin{algorithm}[H]
        \caption{Some more messy pseudocode}\label{algor2}
            \begin{algorithmic}[1]
                \Procedure{SomeProcedure2}{}
                    \State j=0
                    \For{i=0; i$<$n; i++}
                        \State j=i+j
                    \EndFor
                \EndProcedure
            \end{algorithmic}
        \end{algorithm}
    \begin{oneparchoices}
        \choice $\mathcal{O}(\log n)$
        \choice $\mathcal{O}(n\log n)$
        \choice $\mathcal{O}(n^{2})$
        \CorrectChoice $\mathcal{O}(n)$
    \end{oneparchoices}


\question Using the provided pseudo-code, find the worst case performance in Big-$\mathcal{O}$ notation. Assume we know that \textbf{someMethod(n)} is $\mathcal{O}(\log n)$.
        \begin{algorithm}[H]
        \caption{Some pseudocode with a method in it}\label{algor3}
            \begin{algorithmic}[1] % The number tells where the line numbering should start
                \Procedure{SomeProcedure3}{}
                    \State j=0
                    \For{i=0; i$<$n; i++}
                        \State j=someMethod(n)
                    \EndFor
                \EndProcedure
            \end{algorithmic}
        \end{algorithm}
    \begin{oneparchoices}
        \choice $\mathcal{O}(\log n)$
        \CorrectChoice $\mathcal{O}(n\log n)$
        \choice $\mathcal{O}(n^{2})$
        \choice $\mathcal{O}(n)$
    \end{oneparchoices}


\question Using the provided pseudo-code, find the worst case performance in Big-$\mathcal{O}$ notation.
        \begin{algorithm}[H]
        \caption{Some more messy pseudocode}\label{algor4}
            \begin{algorithmic}[1]
                \Procedure{SomeProcedure4}{}
                    \While {n$>1$}
                        \State n=n/2
                    \EndWhile
                \EndProcedure
            \end{algorithmic}
        \end{algorithm}
    \begin{oneparchoices}
        \CorrectChoice $\mathcal{O}(\log n)$
        \choice $\mathcal{O}(n\log n)$
        \choice $\mathcal{O}(n^{2})$
        \choice $\mathcal{O}(n)$
    \end{oneparchoices}


% \question Using the provided pseudo-code, find the worst case performance in Big-$\mathcal{O}$ notation.
%     \begin{algorithm}[H]
%     \caption{Euclid’s algorithm}
%         \label{euclid}
%         \begin{algorithmic}[1] % The number tells where the line numbering should start
%             \Procedure{Euclid}{$a,b$} \Comment{The g.c.d. of a and b}
%                 \State $r\gets a \bmod b$
%                 \While{$r\not=0$} \Comment{We have the answer if r is 0}
%                     \State $a \gets b$
%                     \State $b \gets r$
%                     \State $r \gets a \bmod b$
%                 \EndWhile\label{euclidendwhile}
%                 \State \textbf{return} $b$\Comment{The gcd is b}
%             \EndProcedure
%         \end{algorithmic}
%     \end{algorithm}
%     \begin{oneparchoices}
%         \choice $\mathcal{O}(\log n)$
%         \choice $\mathcal{O}(n\log n)$
%         \choice $\mathcal{O}(n^{2})$
%         \CorrectChoice $\mathcal{O}(n)$
%     \end{oneparchoices}
\end{questions}

\vfill
\subsection*{Key}
C, D, D, B, A

\end{document}